\documentclass[10pt,a4paper]{amsart}
\usepackage{amsmath,amssymb,amsthm,graphicx}
\usepackage{verbatim}
\newenvironment{reference}{\comment}{\endcomment}
%\newenvironment{reference}{}{}
\newenvironment{slogan}{\comment}{\endcomment}
\newenvironment{history}{\comment}{\endcomment}

% For commutative diagrams we use Xy-pic
\usepackage[all]{xy}

% We use 2cell for 2-commutative diagrams.
\xyoption{2cell}
% Theorem environments.
%
\theoremstyle{plain}
\newtheorem{theorem}[subsection]{Theorem}
\newtheorem{proposition}[subsection]{Proposition}
\newtheorem{lemma}[subsection]{Lemma}

\theoremstyle{definition}
\newtheorem{definition}[subsection]{Definition}
\newtheorem{example}[subsection]{Example}
\newtheorem{exercise}[subsection]{Exercise}
\newtheorem{situation}[subsection]{Situation}

\theoremstyle{remark}
\newtheorem{remark}[subsection]{Remark}
\newtheorem{remarks}[subsection]{Remarks}

\numberwithin{equation}{subsection}

% Macros
%
\def\lim{\mathop{\mathrm{lim}}\nolimits}
\def\colim{\mathop{\mathrm{colim}}\nolimits}
\def\Spec{\mathop{\mathrm{Spec}}}
\def\Hom{\mathop{\mathrm{Hom}}\nolimits}
\def\Ext{\mathop{\mathrm{Ext}}\nolimits}
\def\SheafHom{\mathop{\mathcal{H}\!\mathit{om}}\nolimits}
\def\SheafExt{\mathop{\mathcal{E}\!\mathit{xt}}\nolimits}
\def\Sch{\mathit{Sch}}
\def\Mor{\mathop{\mathrm{Mor}}\nolimits}
\def\Ob{\mathop{\mathrm{Ob}}\nolimits}
\def\Sh{\mathop{\mathit{Sh}}\nolimits}
\def\NL{\mathop{N\!L}\nolimits}
\def\CH{\mathop{\mathrm{CH}}\nolimits}
\def\proetale{{pro\text{-}\acute{e}tale}}
\def\etale{{\acute{e}tale}}
\def\QCoh{\mathit{QCoh}}
\def\Ker{\mathop{\mathrm{Ker}}}
\def\Im{\mathop{\mathrm{Im}}}
\def\Coker{\mathop{\mathrm{Coker}}}
\def\Coim{\mathop{\mathrm{Coim}}}

% Boxtimes
%
\DeclareMathSymbol{\boxtimes}{\mathbin}{AMSa}{"02}

%
% Macros for moduli stacks/spaces
%
\def\QCohstack{\mathcal{QC}\!\mathit{oh}}
\def\Cohstack{\mathcal{C}\!\mathit{oh}}
\def\Spacesstack{\mathcal{S}\!\mathit{paces}}
\def\Quotfunctor{\mathrm{Quot}}
\def\Hilbfunctor{\mathrm{Hilb}}
\def\Curvesstack{\mathcal{C}\!\mathit{urves}}
\def\Polarizedstack{\mathcal{P}\!\mathit{olarized}}
\def\Complexesstack{\mathcal{C}\!\mathit{omplexes}}
% \Pic is the operator that assigns to X its picard group, usage \Pic(X)
% \Picardstack_{X/B} denotes the Picard stack of X over B
% \Picardfunctor_{X/B} denotes the Picard functor of X over B
\def\Pic{\mathop{\mathrm{Pic}}\nolimits}
\def\Picardstack{\mathcal{P}\!\mathit{ic}}
\def\Picardfunctor{\mathrm{Pic}}
\def\Deformationcategory{\mathcal{D}\!\mathit{ef}}
\usepackage{cleveref}
\crefname{originalreference}{original reference}{original references}
\Crefname{originalreference}{Original reference}{Original references}
\makeatletter
\newlabel{situation-versal}{{0DQU}{source}{}{sourcepointer-1}{}}
\newlabel{situation-versal@cref}{{[originalreference][0][]0DQU}{[1][0][]source}{}{}{}}
\makeatother
\crefname{originalreference}{original reference}{original references}
\Crefname{originalreference}{Original reference}{Original references}
\makeatletter
\newlabel{stacks-morphisms-lemma-where-smooth}{{0DZR}{source}{}{sourcepointer-2}{}}
\newlabel{stacks-morphisms-lemma-where-smooth@cref}{{[originalreference][0][]0DZR}{[1][0][]source}{}{}{}}
\makeatother
\crefname{originalreference}{original reference}{original references}
\Crefname{originalreference}{Original reference}{Original references}
\makeatletter
\newlabel{formal-defos-lemma-smooth-power-series-classical}{{0DZK}{source}{}{sourcepointer-3}{}}
\newlabel{formal-defos-lemma-smooth-power-series-classical@cref}{{[originalreference][0][]0DZK}{[1][0][]source}{}{}{}}
\makeatother
\crefname{originallemma}{original lemma}{original lemmas}
\Crefname{originallemma}{Original lemma}{Original lemmas}
\makeatletter
\newlabel{lemma-versal-ring}{{0DQX}{source}{}{sourcepointer-4}{}}
\newlabel{lemma-versal-ring@cref}{{[originallemma][0][]0DQX}{[1][0][]source}{}{}{}}
\makeatother
\crefname{originallemma}{original lemma}{original lemmas}
\Crefname{originallemma}{Original lemma}{Original lemmas}
\makeatletter
\newlabel{lemma-versal-ring-field-extension}{{0DQY}{source}{}{sourcepointer-5}{}}
\newlabel{lemma-versal-ring-field-extension@cref}{{[originallemma][0][]0DQY}{[1][0][]source}{}{}{}}
\makeatother
\crefname{originalreference}{original reference}{original references}
\Crefname{originalreference}{Original reference}{Original references}
\makeatletter
\newlabel{more-algebra-lemma-compose-formally-smooth}{{07EF}{source}{}{sourcepointer-6}{}}
\newlabel{more-algebra-lemma-compose-formally-smooth@cref}{{[originalreference][0][]07EF}{[1][0][]source}{}{}{}}
\makeatother
\crefname{originalreference}{original reference}{original references}
\Crefname{originalreference}{Original reference}{Original references}
\makeatletter
\newlabel{more-algebra-proposition-fs-regular}{{07PM}{source}{}{sourcepointer-7}{}}
\newlabel{more-algebra-proposition-fs-regular@cref}{{[originalreference][0][]07PM}{[1][0][]source}{}{}{}}
\makeatother
\crefname{originalreference}{original reference}{original references}
\Crefname{originalreference}{Original reference}{Original references}
\makeatletter
\newlabel{more-algebra-lemma-regular-permanence}{{07NT}{source}{}{sourcepointer-8}{}}
\newlabel{more-algebra-lemma-regular-permanence@cref}{{[originalreference][0][]07NT}{[1][0][]source}{}{}{}}
\makeatother
\crefname{originallemma}{original lemma}{original lemmas}
\Crefname{originallemma}{Original lemma}{Original lemmas}
\makeatletter
\newlabel{lemma-compare-versal-ring-completion}{{0DQZ}{source}{}{sourcepointer-9}{}}
\newlabel{lemma-compare-versal-ring-completion@cref}{{[originallemma][0][]0DQZ}{[1][0][]source}{}{}{}}
\makeatother
\crefname{originalreference}{original reference}{original references}
\Crefname{originalreference}{Original reference}{Original references}
\makeatletter
\newlabel{more-algebra-lemma-formally-smooth-finite-type}{{07VH}{source}{}{sourcepointer-10}{}}
\newlabel{more-algebra-lemma-formally-smooth-finite-type@cref}{{[originalreference][0][]07VH}{[1][0][]source}{}{}{}}
\makeatother
\title{Smooth locus, formal smoothness of a versal ring, and unobstructedness are equivalent}
\author{The Stacks Project Authors}
\date{stacks-project:stacks-geometry.tex:local-snapshot-20261012; GNU FDL 1.2 or later}
\begin{document}
\maketitle
\noindent\textit{Editorial note.} All delivered mathematical text is the author's own, copied verbatim from the pinned Stacks Project chapter file stacks-geometry.tex; no statement or proof step was rewritten, reordered or removed.\par\medskip
\noindent\textit{Editorial note.} Span recomputation by label pairing (LANE\_BRIEF 10.2): in this chapter the author's proof environment follows the statement environment as a standalone block; the spec covers the complete environment pair (statement span, proof span) as recorded in the pool candidate.\par\medskip
\noindent\textit{Editorial note.} Occupation scan: combined claim\_key + statement\_anchor + generic-line scan over all registration files (active and withdrawn) found 0 hits for stacks-geometry.tex (ge-occupied.json, 2026-10-14); 207 route=geometry / chapter-list field hits in focused-seeds are alias noise (route name and alphabetical chapter enumerations), excluded per the ladder alias-blind-spot caveat.\par\medskip
\ifdefined\hypertarget\hypertarget{sourcepointer-1}{}\fi
\noindent\textit{Source reference.} Original reference 0DQU: sources/stacks-chow/tags.txt, source lines 17199--17199. Cross-chapter standard prerequisite cited by the author inside the counted span; it is used as a pointer and its own proof is not part of this item.\par\medskip
\ifdefined\hypertarget\hypertarget{sourcepointer-2}{}\fi
\noindent\textit{Source reference.} Original reference 0DZR: sources/stacks-chow/tags.txt, source lines 17511--17511. Cross-chapter standard prerequisite cited by the author inside the counted span; it is used as a pointer and its own proof is not part of this item.\par\medskip
\ifdefined\hypertarget\hypertarget{sourcepointer-3}{}\fi
\noindent\textit{Source reference.} Original reference 0DZK: sources/stacks-chow/tags.txt, source lines 17505--17505. Cross-chapter standard prerequisite cited by the author inside the counted span; it is used as a pointer and its own proof is not part of this item.\par\medskip
\ifdefined\hypertarget\hypertarget{sourcepointer-4}{}\fi
\noindent\textit{Source reference.} Original lemma 0DQX: sources/stacks-chow/tags.txt, source lines 17202--17202. Cross-chapter standard prerequisite cited by the author inside the counted span; it is used as a pointer and its own proof is not part of this item.\par\medskip
\ifdefined\hypertarget\hypertarget{sourcepointer-5}{}\fi
\noindent\textit{Source reference.} Original lemma 0DQY: sources/stacks-chow/tags.txt, source lines 17203--17203. Cross-chapter standard prerequisite cited by the author inside the counted span; it is used as a pointer and its own proof is not part of this item.\par\medskip
\ifdefined\hypertarget\hypertarget{sourcepointer-6}{}\fi
\noindent\textit{Source reference.} Original reference 07EF: sources/stacks-chow/tags.txt, source lines 9445--9445. Cross-chapter standard prerequisite cited by the author inside the counted span; it is used as a pointer and its own proof is not part of this item.\par\medskip
\ifdefined\hypertarget\hypertarget{sourcepointer-7}{}\fi
\noindent\textit{Source reference.} Original reference 07PM: sources/stacks-chow/tags.txt, source lines 9802--9802. Cross-chapter standard prerequisite cited by the author inside the counted span; it is used as a pointer and its own proof is not part of this item.\par\medskip
\ifdefined\hypertarget\hypertarget{sourcepointer-8}{}\fi
\noindent\textit{Source reference.} Original reference 07NT: sources/stacks-chow/tags.txt, source lines 9773--9773. Cross-chapter standard prerequisite cited by the author inside the counted span; it is used as a pointer and its own proof is not part of this item.\par\medskip
\ifdefined\hypertarget\hypertarget{sourcepointer-9}{}\fi
\noindent\textit{Source reference.} Original lemma 0DQZ: sources/stacks-chow/tags.txt, source lines 17204--17204. Cross-chapter standard prerequisite cited by the author inside the counted span; it is used as a pointer and its own proof is not part of this item.\par\medskip
\ifdefined\hypertarget\hypertarget{sourcepointer-10}{}\fi
\noindent\textit{Source reference.} Original reference 07VH: sources/stacks-chow/tags.txt, source lines 10007--10007. Cross-chapter standard prerequisite cited by the author inside the counted span; it is used as a pointer and its own proof is not part of this item.\par\medskip
\begin{lemma}
\label{lemma-characterize-smoothness}
In Situation \ref{situation-versal}. Let $x_0 : \Spec(k) \to \mathcal{X}$
be a morphism such that $\Spec(k) \to S$ is of finite type with image $s$.
Let $A$ be a versal ring to $\mathcal{X}$ at $x_0$. The following
are equivalent
\begin{enumerate}
\item $x_0$ is in the smooth locus of $\mathcal{X} \to S$
(Morphisms of Stacks, Lemma \ref{stacks-morphisms-lemma-where-smooth}),
\item $\mathcal{O}_{S, s} \to A$ is formally smooth in the
$\mathfrak m_A$-adic topology, and
\item $\mathcal{F}_{\mathcal{X}, k, x_0}$ is unobstructed.
\end{enumerate}
\end{lemma}

\begin{proof}
The equivalence of (2) and (3) follows immediately from
Formal Deformation Theory, Lemma
\ref{formal-defos-lemma-smooth-power-series-classical}.

\medskip\noindent
Note that $\mathcal{O}_{S, s} \to A$ is formally smooth in the
$\mathfrak m_A$-adic topology if and only if
$\mathcal{O}_{S, s} \to A' = A[[t_1, \ldots, t_r]]$
is formally smooth in the $\mathfrak m_{A'}$-adic topology.
Hence (2) does not depend on the choice
of our versal ring by Lemma \ref{lemma-versal-ring}.
Next, let $l/k$ be a finite extension and choose
$A \to A'$ as in Lemma \ref{lemma-versal-ring-field-extension}.
If $\mathcal{O}_{S, s} \to A$ is formally smooth in the
$\mathfrak m_A$-adic topology, then $\mathcal{O}_{S, s} \to A'$
is formally smooth in the $\mathfrak m_{A'}$-adic topology, see
More on Algebra, Lemma \ref{more-algebra-lemma-compose-formally-smooth}.
Conversely, if $\mathcal{O}_{S, s} \to A'$
is formally smooth in the $\mathfrak m_{A'}$-adic topology,
then $\mathcal{O}_{S, s}^\wedge \to A'$ and $A \to A'$ are regular
(More on Algebra, Proposition \ref{more-algebra-proposition-fs-regular})
and hence $\mathcal{O}_{S, s}^\wedge \to A$ is regular
(More on Algebra, Lemma \ref{more-algebra-lemma-regular-permanence}),
hence $\mathcal{O}_{S, s} \to A$ is formally smooth in the
$\mathfrak m_A$-adic topology (same lemma as before).
Thus the equivalence of (2) and (1) holds for
$k$ and $x_0$ if and only if it holds for $l$ and $x_{0, l}$.

\medskip\noindent
Choose a scheme $U$ and a smooth morphism $U \to \mathcal{X}$ such
that $\Spec(k) \times_\mathcal{X} U$ is nonempty. Choose a finite
extension $l/k$ and a point $w_0 : \Spec(l) \to \Spec(k) \times_\mathcal{X} U$.
Let $u_0 \in U$ be the image of $w_0$.
We may apply the above to $l/k$ and to $l/\kappa(u_0)$
to see that we can reduce to $u_0$. Thus we may assume
$A = \mathcal{O}_{U, u_0}^\wedge$, see
Lemma \ref{lemma-compare-versal-ring-completion}.
Observe that $x_0$ is in the smooth locus of $\mathcal{X} \to S$
if and only if $u_0$ is in the smooth locus of $U \to S$, see for example
Morphisms of Stacks, Lemma \ref{stacks-morphisms-lemma-where-smooth}.
Thus the equivalence of (1) and (2) follows from
More on Algebra, Lemma \ref{more-algebra-lemma-formally-smooth-finite-type}.
\end{proof}

\end{document}
